La presentación de resultados del preprocesamiento consiste en describir el número de registros originales y conservados por fuente, temporada y equipo; los valores faltantes; los intervalos y posesiones reconstruidos; y las observaciones descartadas con su justificación. También se documentarán distribuciones de exposición de quintetos y ejemplos de tablas finales de entrenamiento, sin confundirlos con resultados de predicción.

Para evaluar los modelos en una etapa posterior se comunicarán tanto el \textit{Net Rating} esperado como la probabilidad de diferencial positivo, acompañados de métricas de error, calibración y ordenamiento de recomendaciones. La presentación distinguirá resultados experimentales medidos de objetivos metodológicos aún pendientes.